\documentclass[10pt,a4paper]{amsart}
\usepackage{amsmath,amssymb,amsthm,graphicx}
\usepackage{verbatim}
\newenvironment{reference}{\comment}{\endcomment}
%\newenvironment{reference}{}{}
\newenvironment{slogan}{\comment}{\endcomment}
\newenvironment{history}{\comment}{\endcomment}

% For commutative diagrams we use Xy-pic
\usepackage[all]{xy}

% We use 2cell for 2-commutative diagrams.
\xyoption{2cell}
% Theorem environments.
%
\theoremstyle{plain}
\newtheorem{theorem}[subsection]{Theorem}
\newtheorem{proposition}[subsection]{Proposition}
\newtheorem{lemma}[subsection]{Lemma}

\theoremstyle{definition}
\newtheorem{definition}[subsection]{Definition}
\newtheorem{example}[subsection]{Example}
\newtheorem{exercise}[subsection]{Exercise}
\newtheorem{situation}[subsection]{Situation}

\theoremstyle{remark}
\newtheorem{remark}[subsection]{Remark}
\newtheorem{remarks}[subsection]{Remarks}

\numberwithin{equation}{subsection}

% Macros
%
\def\lim{\mathop{\mathrm{lim}}\nolimits}
\def\colim{\mathop{\mathrm{colim}}\nolimits}
\def\Spec{\mathop{\mathrm{Spec}}}
\def\Hom{\mathop{\mathrm{Hom}}\nolimits}
\def\Ext{\mathop{\mathrm{Ext}}\nolimits}
\def\SheafHom{\mathop{\mathcal{H}\!\mathit{om}}\nolimits}
\def\SheafExt{\mathop{\mathcal{E}\!\mathit{xt}}\nolimits}
\def\Sch{\mathit{Sch}}
\def\Mor{\mathop{\mathrm{Mor}}\nolimits}
\def\Ob{\mathop{\mathrm{Ob}}\nolimits}
\def\Sh{\mathop{\mathit{Sh}}\nolimits}
\def\NL{\mathop{N\!L}\nolimits}
\def\CH{\mathop{\mathrm{CH}}\nolimits}
\def\proetale{{pro\text{-}\acute{e}tale}}
\def\etale{{\acute{e}tale}}
\def\QCoh{\mathit{QCoh}}
\def\Ker{\mathop{\mathrm{Ker}}}
\def\Im{\mathop{\mathrm{Im}}}
\def\Coker{\mathop{\mathrm{Coker}}}
\def\Coim{\mathop{\mathrm{Coim}}}

% Boxtimes
%
\DeclareMathSymbol{\boxtimes}{\mathbin}{AMSa}{"02}

%
% Macros for moduli stacks/spaces
%
\def\QCohstack{\mathcal{QC}\!\mathit{oh}}
\def\Cohstack{\mathcal{C}\!\mathit{oh}}
\def\Spacesstack{\mathcal{S}\!\mathit{paces}}
\def\Quotfunctor{\mathrm{Quot}}
\def\Hilbfunctor{\mathrm{Hilb}}
\def\Curvesstack{\mathcal{C}\!\mathit{urves}}
\def\Polarizedstack{\mathcal{P}\!\mathit{olarized}}
\def\Complexesstack{\mathcal{C}\!\mathit{omplexes}}
% \Pic is the operator that assigns to X its picard group, usage \Pic(X)
% \Picardstack_{X/B} denotes the Picard stack of X over B
% \Picardfunctor_{X/B} denotes the Picard functor of X over B
\def\Pic{\mathop{\mathrm{Pic}}\nolimits}
\def\Picardstack{\mathcal{P}\!\mathit{ic}}
\def\Picardfunctor{\mathrm{Pic}}
\def\Deformationcategory{\mathcal{D}\!\mathit{ef}}
\usepackage{cleveref}
\crefname{originallemma}{original lemma}{original lemmas}
\Crefname{originallemma}{Original lemma}{Original lemmas}
\makeatletter
\newlabel{lemma-additive-functor}{{0DLP}{source}{}{sourcepointer-1}{}}
\newlabel{lemma-additive-functor@cref}{{[originallemma][0][]0DLP}{[1][0][]source}{}{}{}}
\makeatother
\crefname{originallemma}{original lemma}{original lemmas}
\Crefname{originallemma}{Original lemma}{Original lemmas}
\makeatletter
\newlabel{lemma-additive-cat-biproduct-kernel}{{09QG}{source}{}{sourcepointer-2}{}}
\newlabel{lemma-additive-cat-biproduct-kernel@cref}{{[originallemma][0][]09QG}{[1][0][]source}{}{}{}}
\makeatother
\crefname{originallemma}{original lemma}{original lemmas}
\Crefname{originallemma}{Original lemma}{Original lemmas}
\makeatletter
\newlabel{lemma-snake}{{010H}{source}{}{sourcepointer-3}{}}
\newlabel{lemma-snake@cref}{{[originallemma][0][]010H}{[1][0][]source}{}{}{}}
\makeatother
\title{Characterizations of additive, left exact, right exact and exact functors}
\author{The Stacks Project Authors}
\date{stacks-project:homology.tex:local-snapshot-20261012; GNU FDL 1.2 or later}
\begin{document}
\maketitle
\noindent\textit{Editorial note.} All delivered mathematical text is the author's own, copied from the pinned Stacks Project chapter file homology.tex; no statement or proof step was rewritten, reordered or removed.\par\medskip
\noindent\textit{Editorial note.} The proof is the author's complete argument for the four statements; the sentence 'The proof of (3) is similar to the proof of (2)' points to the fully written proof of the second statement, of which the third is the exact dual, and it is preserved verbatim inside the counted span.\par\medskip
\noindent\textit{Editorial note.} Every labelled reference occurring inside the counted spans is carried in assembly.external\_references as a non-load-bearing navigation pointer (same-chapter ones at their definition site, cross-chapter ones at the invocation inside the counted span), so no load-bearing mathematics is replaced by a stub. No bibliography citation occurs inside the counted spans.\par\medskip
\noindent\textit{Editorial note.} The same-name check for this card scanned every staged Stacks chapter of this repository -- all sources/stacks-*/pinned-original-source.tex files, the card's own chapter excluded -- and the labels found are recorded in dedup.same\_name\_cross\_chapter on the card.\par\medskip
\noindent\textit{Editorial note.} The preamble is the chapter's own native preamble excerpt already used for this source: the comment-type reference, slogan and history environments (native preamble lines 8-13), the Xy-pic setup (native preamble lines 14-18) and the native theorem declarations with the mathematical macros (native preamble lines 154-217, the actual last line of the file).\par\medskip
\noindent\textit{Allowed standard prerequisites (editorial).}\begin{itemize}
\item Additive functors and preservation of direct sums. Same chapter, Lemma lemma-additive-functor and Lemma lemma-additive-cat-biproduct-kernel. Supplies the characterization of additivity used in the first part and the short exact sequence involving a direct sum.
\end{itemize}\medskip
\ifdefined\hypertarget\hypertarget{sourcepointer-1}{}\fi
\noindent\textit{Source reference.} Original lemma 0DLP: sources/stacks-homology/pinned-original-source.tex, source lines 1687--1687. Navigation pointer invoked at proof line 1786; the counted span contains the pointer only and the cited result is not reproduced inside it, so the pointer is registered as a non-load-bearing navigation stub in the established form for Stacks material. No bibliography citation is involved.\par\medskip
\ifdefined\hypertarget\hypertarget{sourcepointer-2}{}\fi
\noindent\textit{Source reference.} Original lemma 09QG: sources/stacks-homology/pinned-original-source.tex, source lines 470--470. Navigation pointer invoked at proof line 1794; the counted span contains the pointer only and the cited result is not reproduced inside it, so the pointer is registered as a non-load-bearing navigation stub in the established form for Stacks material. No bibliography citation is involved.\par\medskip
\ifdefined\hypertarget\hypertarget{sourcepointer-3}{}\fi
\noindent\textit{Source reference.} Original lemma 010H: sources/stacks-homology/pinned-original-source.tex, source lines 1229--1229. Navigation pointer invoked at proof line 1809; the counted span contains the pointer only and the cited result is not reproduced inside it, so the pointer is registered as a non-load-bearing navigation stub in the established form for Stacks material. No bibliography citation is involved.\par\medskip
\begin{lemma}
\label{lemma-exact-functor}
Let $\mathcal{A}$ and $\mathcal{B}$ be abelian categories.
Let $F : \mathcal{A} \to \mathcal{B}$ be a functor.
\begin{enumerate}
\item If $F$ is either left or right exact, then it is additive.
\item $F$ is left exact if and only if
for every short exact sequence
$0 \to A \to B \to C \to 0$
the sequence $0 \to F(A) \to F(B) \to F(C)$
is exact.
\item $F$ is right exact if and only if for every short exact sequence
$0 \to A \to B \to C \to 0$
the sequence $F(A) \to F(B) \to F(C) \to 0$
is exact.
\item $F$ is exact if and only if for every short exact sequence
$0 \to A \to B \to C \to 0$
the sequence $0 \to F(A) \to F(B) \to F(C) \to 0$
is exact.
\end{enumerate}
\end{lemma}

\begin{proof}
If $F$ is left exact, i.e., $F$ commutes with finite limits, then
$F$ sends products to products, hence $F$ preserved direct sums,
hence $F$ is additive by Lemma \ref{lemma-additive-functor}.
On the other hand, suppose that for every short exact sequence
$0 \to A \to B \to C \to 0$ the sequence $0 \to F(A) \to F(B) \to F(C)$
is exact. Let $A, B$ be two objects. Then we have a short
exact sequence
$$
0 \to A \to A \oplus B \to B \to 0
$$
see for example Lemma \ref{lemma-additive-cat-biproduct-kernel}.
By assumption, the lower row in the commutative diagram
$$
\xymatrix{
0 \ar[r] &
F(A) \ar[d] \ar[r] &
F(A) \oplus F(B) \ar[r] \ar[d] &
F(B) \ar[d] \ar[r] &
0 \\
0 \ar[r] &
F(A) \ar[r] &
F(A \oplus B) \ar[r] &
F(B)
}
$$
is exact. Hence by the snake lemma (Lemma \ref{lemma-snake})
we conclude that $F(A) \oplus F(B) \to F(A \oplus B)$ is an
isomorphism. Hence $F$ is additive in this case as well.
Thus for the rest of the proof we may assume $F$ is additive.

\medskip\noindent
Denote $f : B \to C$ a map from $B$ to $C$.
Exactness of $0 \to A \to B \to C$ just means that
$A = \Ker(f)$. Clearly the kernel of $f$ is
the equalizer of the two maps $f$ and $0$ from $B$ to $C$.
Hence if $F$ commutes with limits, then $F(\Ker(f))
= \Ker(F(f))$ which exactly means that
$0 \to F(A) \to F(B) \to F(C)$ is exact.

\medskip\noindent
Conversely, suppose that $F$ is additive and
transforms any short exact sequence $0 \to A \to B \to C \to 0$ into
an exact sequence $0 \to F(A) \to F(B) \to F(C)$.
Because it is additive it commutes with direct sums
and hence finite products in $\mathcal{A}$. To show
it commutes with finite limits it therefore
suffices to show that it commutes with
equalizers. But equalizers in an abelian category
are the same as the kernel of the difference map,
hence it suffices to show that $F$ commutes with
taking kernels. Let $f : A \to B$ be a morphism.
Factor $f$ as $A \to I \to B$ with $f' : A \to I$ surjective
and $i : I \to B$ injective. (This is possible by the
definition of an abelian category.) Then it is
clear that $\Ker(f) = \Ker(f')$. Also
$0 \to \Ker(f') \to A \to I \to 0$
and
$0 \to I \to B \to B/I \to 0$
are short exact. By the condition imposed on $F$
we see that
$0 \to F(\Ker(f')) \to F(A) \to F(I)$
and
$0 \to F(I) \to F(B) \to F(B/I)$
are exact. Hence it is also the case that
$F(\Ker(f'))$ is the kernel of the map
$F(A) \to F(B)$, and we win.

\medskip\noindent
The proof of (3) is similar to the proof of (2).
Statement (4) is a combination of (2) and (3).
\end{proof}

\end{document}
